\chapter{評価結果と考察}
\label{chap:evaluation}

本章では，第\ref{chap:experiment}章で述べた実験シナリオに基づく評価結果を示し，提案手法の有効性と課題について考察する．

\section{評価結果}

\subsection{C-Plane接続検証}
\label{subsec:cplane}
\paragraph{結果}
sXGP基地局配下のUE（Pixel）から，提案ゲートウェイを経由してOpen5GSへのRegistration（登録）手順を実行した結果，以下の動作を確認した．

\begin{itemize}
    \item \textbf{S1 Setup / NG Setup}: ゲートウェイ起動時にeNBおよびAMFとの接続が確立され，ID変換（Global eNB ID $\leftrightarrow$ Global gNB ID）が正常に行われた．
    \item \textbf{NASメッセージ変換}: UEが送信したAttach Request（4G NAS）がゲートウェイにおいてRegistration Request（5G NAS）に変換され，AMFに到達した．
    \item \textbf{認証・セキュリティ}: 5G-AKA認証が成功し，セキュリティモード手順（Security Mode Command/Complete）が完了した．
    \item \textbf{PDU Session確立}: Registration完了後，Early S1 Setup手順によりeNB側でベアラが確立され，続いてAMF側でPDU Sessionが確立された．
    \item \textbf{IPアドレス付与}: UEに対してOpen5GSから割り当てられたIPアドレスが付与され，アンテナピクトが表示された．
\end{itemize}

また，ゲートウェイ内部ログに基づく計測では，制御面メッセージの変換処理時間は1 ms未満であった．これはLTE網の登録/アタッチ手順（数秒）と比較して十分に小さく，ユーザ体験に影響を与えない水準である．

\paragraph{定性的評価（設計妥当性の検証）}
本研究の主目的は，提案した変換設計（マッピング規則）の妥当性を検証することである．以下に，pcapおよびログから確認した具体的な観測事実を示す．

\begin{itemize}
    \item \textbf{NAS種別の対応（表\ref{tab:nas_conversion_ul}，表\ref{tab:nas_conversion_dl}の検証）}:
    pcap上で，eNB側の\texttt{Attach Request}（NAS EPS）に対応して，AMF側では\texttt{Registration Request}（NAS 5GS）が観測された．同様に，\texttt{Attach Accept}は\texttt{Registration Accept}として，\texttt{Security Mode Command/Complete}は同名のメッセージとして正しく変換された．

    \item \textbf{ID束の整合（表\ref{tab:id_conversion}の検証）}:
    ゲートウェイログにおいて，eNB側の(ENB-UE-S1AP-ID, MME-UE-S1AP-ID)と5GC側の(RAN-UE-NGAP-ID, AMF-UE-NGAP-ID)が同一UEコンテキストに紐づいて管理されていることを確認した．具体的には，単一のUE接続に対して4つのIDが一貫して対応付けられ，コンテキスト解放まで維持された．

    \item \textbf{AAA機能の成立}:
    Open5GSのログにおいて，認証ベクトル（RAND, AUTN, XRES*）の生成・検証が成功し，\texttt{5GMM state: REGISTERED}への遷移が確認された．これは，4G UEからの認証情報が5G-AKA手順として正しく処理されたことを示す．

    \item \textbf{TEIDマッピングの成立（表\ref{tab:teid_mapping}の検証）}:
    ゲートウェイログにおいて，上り方向のGTP-UパケットでS1-U TEID（eNB割当）がN3 TEID（UPF割当）へ書き換えられていることを確認した．下り方向は透過転送（Hybrid Proxy設計）により，TEIDの書き換えなしで正常に転送された．
\end{itemize}

以上の観測結果は，第\ref{chap:proposal}章で定義した変換規則（NAS変換，ID変換，TEIDマッピング，QoSマッピング）が実際の信号レベルで成立していることを示している．これは，「世代共通機能」という抽象化に基づく設計アプローチの妥当性を裏付けるものである．

\subsection{U-Planeデータ転送性能の評価}
\paragraph{遅延}
UEからコアネットワーク内のサーバへの往復遅延（RTT）を\texttt{irtt}（Isochronous Round-Trip Tester）を用いて測定した．\texttt{irtt}はUDPベースの高精度遅延測定ツールであり，ICMPベースのpingと比較してより正確なネットワーク遅延特性を把握できる．測定対象のサーバはUPFと同一ホスト上に配置し，ネットワーク経路上の外部要因を排除した．
図\ref{fig:rtt_comparison}に，提案手法とベースライン環境におけるRTTの比較を示す．

\begin{figure}[htbp]
  \centering
  \begin{tikzpicture}
    \begin{axis}[
      ybar,
      bar width=1.2cm,
      width=10cm,
      height=6cm,
      ylabel={RTT (ms)},
      symbolic x coords={提案手法, ベースライン},
      xtick=data,
      ymin=0,
      ymax=100,
      nodes near coords,
      every node near coord/.append style={font=\small},
      enlarge x limits=0.5,
    ]
    \addplot coordinates {(提案手法, 81.1) (ベースライン, 77.7)};
    \end{axis}
  \end{tikzpicture}
  \caption{RTT（往復遅延）の比較（irtt 100回測定の平均値）}
  \label{fig:rtt_comparison}
\end{figure}

\begin{itemize}
    \item \textbf{ゲートウェイ経由（提案手法）}: 平均 81.1 ms（標準偏差 4.8 ms）
    \item \textbf{LTEネイティブ（ベースライン）}: 平均 77.7 ms（標準偏差 16.7 ms）
\end{itemize}
両環境のRTT差は約3.4 msであり，LTE無線区間の遅延（約70〜80 ms）に対して十分に小さい．提案手法の標準偏差がベースラインより小さいのは，測定時の無線環境の違いによるものと考えられる．この結果は，ゲートウェイの変換処理（GTP-Uヘッダ書き換え）による追加遅延が実用上無視できるレベルであることを示している．

加えて，ゲートウェイ内部ログに基づく計測では，ユーザ面のGTP-Uヘッダ書き換え処理は0.1 ms未満であった．

\paragraph{スループット}
iPerf3を用いたUDPスループット測定を行った（100秒間転送）．UDPを選択した理由は，TCPの輻輳制御アルゴリズムの影響を排除し，純粋なU-Plane転送性能を評価するためである．
下り方向（DL）は5, 10, 15, 16, 20 Mbps，上り方向（UL）は3, 5, 10 Mbpsの帯域設定で測定を実施した．
図\ref{fig:throughput_dl}および図\ref{fig:throughput_ul}に，提案手法とベースライン環境におけるスループットの比較を示す．

\begin{figure}[htbp]
  \centering
  \begin{tikzpicture}
    \begin{axis}[
      width=12cm,
      height=7cm,
      xlabel={設定帯域 (Mbps)},
      ylabel={測定スループット (Mbps)},
      xmin=0, xmax=22,
      ymin=0, ymax=22,
      xtick={5, 10, 15, 16, 20},
      ytick={0, 5, 10, 15, 20},
      legend style={at={(0.02,0.98)}, anchor=north west},
      grid=major,
      mark size=3pt,
    ]
    % 理想線（設定帯域=測定スループット）
    \addplot[dashed, thick, black!50, domain=0:22] {x};
    \addlegendentry{理想値}
    % 提案手法（GW経由）
    \addplot[mark=square*, thick, GoogleBlue] coordinates {
      (5, 5.00) (10, 10.00) (15, 15.00) (16, 15.57) (20, 16.43)
    };
    \addlegendentry{提案手法}
    % ベースライン（4Gネイティブ）
    \addplot[mark=triangle*, thick, GoogleRed] coordinates {
      (5, 5.00) (10, 10.00) (15, 15.00) (16, 16.00) (20, 14.42)
    };
    \addlegendentry{ベースライン}
    \end{axis}
  \end{tikzpicture}
  \caption{下り方向（DL）UDPスループットの比較}
  \label{fig:throughput_dl}
\end{figure}

\begin{figure}[htbp]
  \centering
  \begin{tikzpicture}
    \begin{axis}[
      width=12cm,
      height=7cm,
      xlabel={設定帯域 (Mbps)},
      ylabel={測定スループット (Mbps)},
      xmin=0, xmax=12,
      ymin=0, ymax=12,
      xtick={3, 5, 10},
      ytick={0, 2, 4, 6, 8, 10},
      legend style={at={(0.02,0.98)}, anchor=north west},
      grid=major,
      mark size=3pt,
    ]
    % 理想線（設定帯域=測定スループット）
    \addplot[dashed, thick, black!50, domain=0:12] {x};
    \addlegendentry{理想値}
    % 提案手法（GW経由）
    \addplot[mark=square*, thick, GoogleBlue] coordinates {
      (3, 3.00) (5, 4.31) (10, 4.31)
    };
    \addlegendentry{提案手法}
    % ベースライン（4Gネイティブ）
    \addplot[mark=triangle*, thick, GoogleRed] coordinates {
      (3, 3.00) (5, 4.29) (10, 4.28)
    };
    \addlegendentry{ベースライン}
    \end{axis}
  \end{tikzpicture}
  \caption{上り方向（UL）UDPスループットの比較}
  \label{fig:throughput_ul}
\end{figure}

\begin{itemize}
    \item \textbf{下り（Downlink）}: 15 Mbpsまでは両環境とも設定帯域を達成（パケットロス 0\%）．16 Mbps以上では提案手法が15.6〜16.4 Mbps，ベースラインが16.0〜14.4 Mbpsとなり，両環境とも飽和傾向を示すが有意な差はない．
    \item \textbf{上り（Uplink）}: 3 Mbpsでは両環境とも設定帯域を達成．5 Mbps以上では両環境とも約4.3 Mbpsで飽和し，有意な差はない．
\end{itemize}
下り方向において15 Mbpsまで両環境で差がないことから，通常の利用帯域ではゲートウェイ処理がボトルネックにならないことが確認された．16 Mbps以上では両環境とも飽和傾向を示すが，これはLTE無線区間の下り方向スケジューリング特性や実験環境固有の要因に起因するものと考えられる．上り方向のスループットが約4.3 Mbpsで飽和するのは，LTE無線区間のアップリンクスケジューリング特性（Grant遅延，バッファ管理等）に起因するものであり，両環境で同様の傾向を示すことからゲートウェイ処理とは無関係である．

重要な点は，実用的な帯域範囲（下り15 Mbps以下，上り3 Mbps）において，提案手法とベースライン環境の間で有意な性能差が見られないことである．この範囲ではパケットロス率も0\%であり，ゲートウェイのGTP-Uヘッダ変換処理が通信のボトルネックとなっていないことが確認された．

\section{考察}

\subsection{設計妥当性の総括}
前節の実験結果に基づき，提案した変換設計の妥当性を総括する．

\paragraph{C-Plane設計の妥当性}
定性的評価（第\ref{subsec:cplane}項）において，NAS変換・ID変換・TEIDマッピングがpcapおよびログ上で成立していることを確認した．これは，第\ref{chap:proposal}章で定義した「世代共通機能」という抽象化が，実際のプロトコル変換において有効に機能することを示している．特に，以下の点が重要である．
\begin{itemize}
    \item AAA機能（4G NAS認証を5G-AKAとして処理）が成立し，Open5GSの\texttt{REGISTERED}状態まで遷移した．
    \item ID束（S1AP IDとNGAP ID）が単一UEコンテキストで一貫して管理された．
    \item TEIDマッピング（Hybrid Proxy設計）が意図通りに動作した．
\end{itemize}

\paragraph{U-Plane設計の妥当性}
RTTおよびUDPスループットの測定結果から，実用的な帯域範囲において提案手法とベースライン環境の間に有意な性能差がないことを確認した．特にUDP測定では，下り方向で15 Mbpsまで設定帯域を100\%達成し，パケットロスも0\%であった．これは，ゲートウェイのGTP-Uヘッダ書き換え処理が通常の利用帯域において通信のボトルネックとなっていないことを示している．したがって，Hybrid Proxy設計（上り方向のTEID変換＋下り方向の透過転送）は，性能オーバーヘッドを抑えつつ機能要件を満たす設計として妥当であったと言える．

\paragraph{実験の限界}
本実験は単一UEでの基本動作検証であり，以下の点は未検証である．
\begin{itemize}
    \item 複数UE同時接続時のスケーラビリティ
    \item 高遅延・高損失環境での安定性
    \item 移動管理機能の完全な検証（初期登録は検証済みだが，TAU・ハンドオーバー等のエリア移動は未検証）
    \item ネットワークスライシング等の5GC固有機能
\end{itemize}
したがって，本実験結果は提案設計の「基本動作の成立」を確認したものであり，「実用性の実証」にはさらなる検証が必要である．

\subsection{疎結合アーキテクチャの意義}
本実験により，RANとCNの物理的な結合を解き，論理的なゲートウェイを介して接続する「疎結合アーキテクチャ」の基本動作が確認された．
従来のアーキテクチャでは，RANとCNが同一世代のプロトコルで密に結合しており，CNを次世代へ移行する際にはRANの更新も同時に求められることが多かった．本手法を用いれば，RANを据え置いたまま，CNのみを5GCへ更新する構成が技術的に可能であることが示された．

加えて，本研究の「疎結合」は単なる互換接続（レガシー機器の延命）に留まらず，運用や機能拡張の観点で価値を持つ．ゲートウェイを境界としてRAN側／CN側の責務が分離されることで，(1) CNの更新頻度（ソフトウェアリリースサイクル）をRANの更新制約から切り離せる，(2) プロトコル差分を局所化できるため，差分吸収ロジックの検証範囲を境界内に閉じ込められる，(3) 境界に観測点を集約できるため，障害解析やトレースが容易になる，といった効果が期待できる．特に(2)は，本論文が設計評価（マッピング成立）を重視する立場において重要であり，変換規則（第\ref{chap:proposal}章）と観測（pcap/ログ）の対応付けが明確になる．

さらに，SCTP終端・変換・中継を担う境界要素を明示的に導入することは，ソフトウェアアーキテクチャにおける「互換層（アダプタ）」に相当する．この互換層を介して，レガシーRANの存在を前提にしつつも，CN側は5GCとしての設計原則（サービス分割，独立スケール，段階的更新）を保ったまま進化できる．

\subsection{世代間互換性の実現アプローチ}
本研究では，既存のプロトコル仕様（S1AP/NGAP，NAS EPS/NAS 5GS）を前提とし，ゲートウェイによる変換でプロトコル差分を吸収するアプローチを採用した．しかし，世代間互換性を実現する方法はこれに限らない．本項では，代替アプローチとして「プロトコル設計自体で世代間互換性を確保する」考え方を検討し，本研究の位置づけを明確にする．

\subsubsection{本研究のアプローチ：変換による差分吸収}
本研究が採用したアプローチは，4Gと5Gの既存プロトコル仕様を変更せず，両者の間にプロトコル変換層（ゲートウェイ）を挿入するものである．第\ref{chap:proposal}章で示した通り，世代不変な機能要件（AAA，Mobility Management，U-Plane）に着目し，これらの機能を実現するプロトコル要素間の対応関係（NAS変換，ID変換，TEIDマッピング等）を定義することで，異なる世代間の相互運用を実現した．

このアプローチの利点は，(1) 既存のeNBおよび5GCに改修が不要であり即時適用可能である，(2) 標準化プロセスを経ずに実装・展開できる，(3) 変換ロジックをゲートウェイに局所化できるため検証範囲が明確である，という点にある．一方，(1) 変換処理に伴う遅延・オーバーヘッドが発生する，(2) 世代ごとに変換ルールの設計・実装が必要となる，(3) 5GC固有機能のフル活用には追加の変換ロジックが必要となる，という課題も存在する．

\subsubsection{代替アプローチ：プロトコル設計による互換性確保}
別のアプローチとして，プロトコル仕様自体を世代間互換性を考慮して設計する方法が考えられる．具体的には，以下のような設計原則に基づくプロトコルを構想できる．

\begin{itemize}
    \item \textbf{MUST（必須）フィールド}: 世代不変な機能要件（AAA，Mobility Management，U-Plane）を実現するために必要なパラメータを，すべての世代で共通の形式・エンコーディングで定義する．例えば，ユーザ識別子，認証ベクトル，セッション識別子，QoSパラメータの基本セット等が該当する．
    \item \textbf{OPTIONAL（任意）フィールド}: 特定世代の固有機能（5Gのスライシング，QoS Flow，6G以降の新機能等）は，拡張フィールドとして定義する．旧世代の機器はこれらのフィールドを無視し，新世代の機器のみが解釈・処理する．
\end{itemize}

このアプローチが実現すれば，(1) 変換処理が不要となりオーバーヘッドがゼロになる，(2) 新世代の機能追加が後方互換性を維持したまま可能になる，(3) 機器間の相互運用性が標準レベルで保証される，という利点が得られる．

\subsubsection{両アプローチの比較と本研究の位置づけ}
表\ref{tab:approach_comparison}に両アプローチの比較を示す．

\begin{table}[htbp]
\centering
\caption{世代間互換性実現アプローチの比較}
\label{tab:approach_comparison}
\begin{tabular}{|p{3cm}|p{5.5cm}|p{5.5cm}|}
\hline
\textbf{観点} & \textbf{変換アプローチ（本研究）} & \textbf{プロトコル設計アプローチ} \\
\hline
適用可能性 & 既存システムに即時適用可能 & 新規プロトコル策定が必要 \\
\hline
標準化 & 不要（実装レベルで対応） & 3GPP等での標準化が必須 \\
\hline
性能オーバーヘッド & 変換処理による遅延あり & なし \\
\hline
拡張性 & 世代ごとに変換ルール追加 & OPTIONALフィールド追加で対応 \\
\hline
既存機器の扱い & そのまま利用可能 & 新プロトコル対応が必要 \\
\hline
\end{tabular}
\end{table}

プロトコル設計アプローチは理想的ではあるが，現実には以下の課題がある．(1) 3GPPにおける標準化プロセスは数年単位の時間を要する，(2) 既存の膨大な4G/5G機器との後方互換性をどう確保するかが難題である，(3) 将来世代（6G以降）で必要となる機能を現時点で予測し設計に織り込むことは困難である．

本研究の変換アプローチは，これらの現実的制約を回避し，既存システムに対して即時に適用可能な解決策を提供する．同時に，本研究で行った「世代不変機能の特定」と「プロトコル要素間の対応関係の分析」（第\ref{chap:proposal}章）は，将来のプロトコル設計においてMUST/OPTIONALフィールドを定義する際の基礎的知見としても活用できる．すなわち，本研究は短期的には変換ゲートウェイとして価値を持ち，長期的にはプロトコル設計への示唆を与えるものである．

\subsection{5GC固有機能を前提とした設計価値の拡張}
本研究の実験は基本動作（登録と疎通）に焦点を当てたが，5GCを採用する動機は「接続できる」ことに加え，5GC固有の機能を利用できる点にある．疎結合アーキテクチャによりレガシーRANを5GCへ収容できれば，RAN側が旧世代であっても，CN側では5GCの運用・制御機能を前提にした設計が可能になる．ここでは「5GC固有機能が使える」前提で，疎結合の価値を整理する．

\paragraph{ネットワークスライシング（運用単位の分離）}
RAN側がスライスを直接選択できない場合でも，CN側では加入者属性，DNN/APN相当情報，あるいは外部ポリシに基づき，スライス（S-NSSAI）やUPF経路を選択する運用が考えられる．無線区間の厳密な資源分離は難しい一方で，コア側での論理分離（ポリシ，課金，監視，経路制御）を導入できるため，特定用途（IoT/閉域/研究用）を他用途から分離して運用する効果は得られる．

\paragraph{ポリシ制御とQoSの高度化}
本プロトタイプではQoSは簡易設定に留まるが，5GC側でPCF等のポリシ制御を前提とすれば，アプリケーション種別に応じたセッション制御や，UPFでのシェーピング・フィルタリングを組み合わせたサービス制御が可能となる．RAN側のQCI表現が粗い場合でも，CN側で「どのセッションをどの経路・帯域で扱うか」という設計自由度が増す点は，疎結合の価値である．

\paragraph{可観測性（Observability）と運用自動化}
5GCはクラウドネイティブ運用（ログ/メトリクス/トレース）と親和性が高い．ゲートウェイを境界として導入すると，(1) 変換の入出力イベント，(2) UEコンテキストの状態，(3) TEIDマッピング更新，といった「世代差分に由来する内部状態」を，運用メトリクスとして集約しやすい．これにより，従来は基地局ベンダ依存になりがちな障害切り分けを，コア側運用として標準化しやすくなる．

\paragraph{マイグレーション戦略（段階導入）の具体化}
疎結合アーキテクチャは，CN先行移行のための互換層として機能するため，段階的な導入戦略が設計として描きやすい．具体的なマイグレーション戦略については，第\ref{subsubsec:migration}項で詳述する．

\subsection{デプロイメント戦略}
本ゲートウェイはソフトウェア（コンテナ）として実装されており，ネットワーク上の任意の場所に配置可能である．配置の粒度として，以下の3パターンが考えられる．
\begin{itemize}
    \item \textbf{(A) RANサイト配置}（基地局と同一拠点）: 各eNBの近傍にゲートウェイを配置する．eNB〜ゲートウェイ間の遅延は最小化されるが，拠点数に応じたインスタンス管理が必要となり，運用負荷が高い．
    \item \textbf{(B) TN内配置}（地域集約局等）: 複数のeNBを集約するポイントにゲートウェイを配置する．運用性とレイテンシのバランスが取れる中間的な選択肢である．
    \item \textbf{(C) コア側配置}（データセンター/クラウド）: 5GCと同一拠点にゲートウェイを配置する．集中管理が可能で運用性が最も高いが，eNB〜ゲートウェイ間の遅延が大きくなる可能性がある．
\end{itemize}

配置選択の主な判断軸は運用性（保守・監視のしやすさ，スケーラビリティ，現地作業の有無）である．
多くのユースケースでは，コア側配置(C)が最も現実的な選択肢となる．理由として，(1) 集中管理により運用コストを最小化できる，(2) eNB側への現地作業が不要となる，(3) クラウド/仮想化基盤上でスケールアウトが容易である，といった点が挙げられる．

一方，eNB〜ゲートウェイ間のバックホール帯域が逼迫している場合や，特定区間で細かなトラフィック制御が必要な場合には，TN内配置(B)やRANサイト配置(A)が有効となる場合もある．
本アーキテクチャはソフトウェア実装であるため，要件の変化に応じて配置を柔軟に変更できる点も強みである．

\paragraph{5GC機能階層と配置判断の関係}
第\ref{subsec:5gc_layered_design_framework}項で導入した5GC機能の3層モデル（基盤層・制御メカニズム層・応用サービス層）に基づき，各層がゲートウェイ配置にどのような影響を与えるかを分析する．表\ref{tab:layer_deployment}に各層の特性と配置への影響を示す．

\begin{table}[htbp]
\centering
\caption{5GC機能階層とゲートウェイ配置の関係}
\label{tab:layer_deployment}
\begin{tabular}{|p{2.5cm}|p{3cm}|p{3cm}|p{4.5cm}|}
\hline
\textbf{層} & \textbf{対応方針} & \textbf{連携先NF} & \textbf{配置への影響} \\
\hline
基盤層 & 隠蔽 & なし（CN内部） & 配置に制約なし \\
\hline
制御メカニズム層 & 変換・整合 & AMF, SMF, PCF & コア側配置が有利 \\
\hline
応用サービス層 & スタブ・段階拡張 & NEF, 外部API & 拡張時に再検討 \\
\hline
\end{tabular}
\end{table}

\begin{itemize}
    \item \textbf{基盤層（SBA，CUPS，NFV）}: 第\ref{subsec:5gc_layered_design_framework}項で述べた通り，この層は提案アーキテクチャから完全に隠蔽される．CN内部のサービスベースアーキテクチャ（SBI通信）はゲートウェイから見えないため，基盤層の実装詳細は配置判断に影響しない．

    \item \textbf{制御メカニズム層（スライス選択，QoS制御）}: ゲートウェイが積極的に変換・整合を行う層である．スライス選択（S-NSSAI付与）はNSSF/AMFとの連携が必要であり，QoSマッピング（QCI→5QI変換）はPCF/SMFとのポリシー連携が前提となる．これらのNFは一般にコア側に配置されるため，ゲートウェイもコア側に配置することでNF間通信の遅延を最小化できる．特に，動的なポリシー変更やスライス切替を伴う運用では，この遅延差が顕著になる．

    \item \textbf{応用サービス層（Network Exposure，5G LAN等）}: 現時点ではスタブまたは将来拡張として位置づけられる層である．NEF（Network Exposure Function）を介した外部API連携や，5G LANサービスの提供を行う場合には，連携先システムとの距離が新たな判断軸となる．ただし，本PoCではこの層は実装対象外であり，将来の機能拡張時に配置を再検討する余地がある．
\end{itemize}

\paragraph{変換処理特性と配置の整合性}
ゲートウェイが担う変換処理の特性も配置判断に影響する．

\begin{itemize}
    \item \textbf{C-Plane変換（NAS変換，ID変換）}: 計算負荷が軽く，頻度も接続時に限定されるため，配置による性能差は小さい．ただし，認証情報（AUSF連携）やセッション情報（SMF連携）の取得が必要なため，コア側配置のほうがラウンドトリップが少ない．

    \item \textbf{U-Plane変換（TEIDマッピング）}: パケットごとにヘッダ書き換えが発生するため，リアルタイム性が求められる．ただし，本PoCで確認した通り処理時間は0.1 ms未満であり，配置による遅延差（数ms〜数十ms）のほうが支配的である．高遅延環境（NTN等）では分散配置によりSCTP区間を分離する意義がある．

    \item \textbf{QoSマッピング}: 静的マッピング（QCI→5QIの固定変換）であれば配置に依存しないが，PCF連携による動的ポリシー適用を行う場合はコア側配置が有利である．
\end{itemize}

\paragraph{論理設計に基づく分離配置の可能性}
本PoCでは単一プロセスとして実装したが，第\ref{chap:proposal}章で示した論理設計に基づけば，C-Plane変換機能とU-Plane変換機能を分離配置する構成も可能である．この設計思想は，5GCにおけるCUPS（Control and User Plane Separation）と同様のアプローチである．

\begin{itemize}
    \item \textbf{C-Plane/U-Plane分離配置}: 第\ref{chap:proposal}章の変換設計において，C-Plane変換（NAS変換，ID変換，SCTP/NGAP変換）とU-Plane変換（TEIDマッピング，GTP-Uヘッダ書き換え）は論理的に独立している．両者が共有するのはUEコンテキスト（ID束，TEID対応表）のみであり，この状態情報を外部化すれば，C-PlaneゲートウェイとU-Planeゲートウェイを物理的に分離できる．具体的には，C-Planeをコア側に配置してNF連携を最適化しつつ，U-Planeをエッジ側に配置してデータパスの遅延を最小化する構成が考えられる．

    \item \textbf{機能ごとの独立スケーリング}: C-Plane変換の負荷はセッション数（同時接続UE数）に比例し，状態管理（UEコンテキスト保持）が必要である．一方，U-Plane変換の負荷はスループット（パケット数）に比例し，TEIDテーブル参照以外はステートレスに近い処理である．この特性の違いを活かし，C-PlaneとU-Planeを独立にスケーリングすることで，効率的なリソース配分が可能となる．例えば，IoTシナリオ（多数UE・低トラフィック）ではC-Planeを重点的にスケールし，eMBBシナリオ（少数UE・高トラフィック）ではU-Planeを重点的にスケールする運用が考えられる．
\end{itemize}

この分離配置は，第\ref{subsec:architecture_requirements}項で定義したアーキテクチャ要件との関係で以下のように整理できる．

\begin{itemize}
    \item \textbf{R1（透過的介在）}: S1インタフェースとN2/N3インタフェースの両端を終端する要件は，分離配置でも満たせる．C-PlaneゲートウェイがS1-MME/N2を終端し，U-PlaneゲートウェイがS1-U/N3を終端する構成となる．
    \item \textbf{R2（プロトコル相互変換）}: 変換ロジック自体は配置に依存しない．C-Plane変換とU-Plane変換が論理的に独立している設計により，分離配置が可能となっている．
    \item \textbf{R3（5GC固有機能の整合）}: スライス選択やQoS制御などNF連携が必要な機能は，C-Planeゲートウェイがコア側に配置されることで効率的に実現できる．
\end{itemize}

ただし，本PoCでは分離配置の実装・検証は行っておらず，UEコンテキストの外部化（共有ストレージやメッセージング）に伴う複雑性・遅延の増加については今後の検討課題である．

\paragraph{配置判断のまとめ}
以上の分析から，5GC機能を活用する運用においてはコア側配置(C)が最適と考えられる．特に，制御メカニズム層の機能（スライス選択，QoS制御，ポリシー連携）を活用する場合，ゲートウェイとコアNF間の通信遅延を最小化することが重要である．

ただし，高遅延環境やエッジ処理など特殊要件がある場合には，分散配置が有効となるケースもある．例えば，エッジMEC（ローカルブレイクアウト）を極低遅延で提供したい場合には，UPFだけでなくゲートウェイもエッジ拠点に配置する構成が考えられる．なお，本アーキテクチャの応用シナリオについては\ref{subsec:application}項で詳述する．

総じて，5GC機能を前提とした運用設計においては，コア側配置を基本としつつ，特殊要件がある場合に限り分散配置を検討する，という方針が妥当であると考えられる．

\subsection{応用可能性}
\label{subsec:application}
本研究で提案した疎結合アーキテクチャは，基本的なプロトコル変換の実証を主目的としているが，設計特性から導かれる応用可能性は広い．なお，以下の応用シナリオは本実験では検証しておらず，設計の汎用性から導かれる可能性として整理する．

\subsubsection{研究開発における低コスト検証環境}
高価なgNBを用いずに，安価なsXGP基地局とOSS 5GCを組み合わせた低コストな5GC検証環境を構築できる．免許不要のsXGP帯域（1.9GHz帯）を利用することで，大学や研究機関が免許申請なしに実端末での5GC機能検証を行うことが可能となる．特に，5GC固有機能（スライシング，ポリシー制御等）の研究開発において，RAN側の制約を受けずにCN側の実装・評価に集中できる環境を提供できる．

\subsubsection{商用導入におけるマイグレーション戦略}
\label{subsubsec:migration}
eNBのハードウェア改修なしに5GCへ接続可能であることから，「CN先行移行」戦略の技術的実現可能性を示す．具体的には，(1) 一部エリア・一部eNBから段階的に5GCへ収容する，(2) ゲートウェイのスケールアウトに合わせて収容範囲を拡大する，(3) 将来的にRAN更新（ng-eNB/gNB化）が進んだ段階で互換層の役割を縮退させる，といった段階的導入が可能となる．この手法は，(1) 既存RANへの投資を保護しつつ5GC機能を利用開始できる，(2) RAN更新とCN更新を独立したスケジュールで計画できる，(3) 移行リスクを分散できる，といった利点を持つ．

\subsubsection{NTN等の物理的制約環境}
衛星やHAPSに搭載されたレガシーRANを改修することなく，地上側のソフトウェア更新のみで次世代CNへ収容可能である．

以下では，NTNへの適用可能性について技術的な背景を補足する．

\paragraph{NTNアーキテクチャと本提案の適用}
非地上系ネットワーク（Non-Terrestrial Network: NTN）では，衛星や成層圏プラットフォーム（High Altitude Platform Station: HAPS）に搭載された基地局が広域カバレッジを提供する．

図\ref{fig:ntn_architecture}にNTNの概念構成を示す．

\begin{figure}[htbp]
  \centering
  \begin{tikzpicture}[scale=0.75, every node/.style={font=\small},
    platform/.style={draw, fill=GoogleYellow!50},
    ground/.style={draw, rounded corners, minimum width=1.5cm, minimum height=0.6cm}]

    % 地上局・CN
    \node[ground, fill=GoogleBlue!40] (gw) at (0,1.5) {地上局};
    \node[ground, fill=GoogleGreen!40] (core) at (0,0) {5GC};
    \draw[<->, thick] (gw) -- (core);

    % 衛星
    \node[platform, circle, minimum size=1.2cm] (sat) at (5,4) {衛星};
    \node[above, font=\scriptsize] at (5,4.8) {(eNB/gNB)};

    % HAPS
    \node[platform, ellipse, minimum width=1.6cm, minimum height=0.6cm] (haps) at (10,3) {HAPS};
    \node[above, font=\scriptsize] at (10,3.5) {(eNB/gNB)};

    % UE
    \node[ground, fill=GoogleLightGray] (ue1) at (5,0) {UE};
    \node[ground, fill=GoogleLightGray] (ue2) at (10,0) {UE};

    % 接続線
    \draw[<->, thick, dashed] (sat) -- (gw) node[midway, above, sloped, font=\scriptsize] {Feeder};
    \draw[<->, thick] (sat) -- (ue1) node[midway, right, font=\scriptsize] {Service};
    \draw[<->, thick, dashed] (haps) -- (gw);
    \draw[<->, thick] (haps) -- (ue2);

  \end{tikzpicture}
  \caption{NTN（非地上系ネットワーク）の概念構成}
  \label{fig:ntn_architecture}
\end{figure}

NTNのアーキテクチャには，衛星が単なる中継器（ベントパイプ）として動作する「トランスペアレント（Transparent）型」と，衛星上で基地局機能を実行する「再生（Regenerative）型」が存在する．
図\ref{fig:ntn_payload_comparison}に両アーキテクチャの比較を示す．

\begin{figure}[htbp]
  \centering
  \begin{tikzpicture}[scale=0.8, every node/.style={font=\small},
    sat/.style={draw, circle, minimum size=1.4cm},
    ground/.style={draw, rounded corners, minimum width=1.6cm, minimum height=0.7cm},
    core/.style={draw, rounded corners, minimum width=1.4cm, minimum height=0.6cm, fill=GoogleGreen!40}]

    % ========== Transparent型（左側）==========
    \node[font=\bfseries] at (2.5,5.5) {Transparent型};

    \node[sat, fill=GoogleLightGray!40] (sat-t) at (2.5,3.8) {衛星};
    \node[ground, fill=GoogleBlue!30] (gw-t) at (0.5,1.8) {gNB};
    \node[core] (core-t) at (0.5,0.3) {5GC};
    \node[ground, fill=GoogleLightGray] (ue-t) at (4.5,0.3) {UE};

    \draw[<->, thick] (sat-t) -- (gw-t) node[midway, left, font=\scriptsize] {Feeder};
    \draw[<->, thick] (sat-t) -- (ue-t) node[midway, right, font=\scriptsize] {Service};
    \draw[<->, thick] (gw-t) -- (core-t);

    % ========== 区切り線 ==========
    \draw[dashed, gray] (6.5,6) -- (6.5,-0.5);

    % ========== Regenerative型（右側）==========
    \node[font=\bfseries] at (10,5.5) {Regenerative型};

    \node[sat, fill=GoogleYellow!50] (sat-r) at (10,3.8) {gNB};
    \node[ground, fill=GoogleBlue!30] (gw-r) at (8,1.8) {GW};
    \node[core] (core-r) at (8,0.3) {5GC};
    \node[ground, fill=GoogleLightGray] (ue-r) at (12,0.3) {UE};

    \draw[<->, thick] (sat-r) -- (gw-r) node[midway, left, font=\scriptsize] {Feeder};
    \draw[<->, thick] (sat-r) -- (ue-r) node[midway, right, font=\scriptsize] {Service};
    \draw[<->, thick] (gw-r) -- (core-r);

  \end{tikzpicture}
  \caption{NTNペイロードアーキテクチャの比較（TR 38.821\cite{3GPP_TR38821}に基づく）}
  \label{fig:ntn_payload_comparison}
\end{figure}

Transparent型では，衛星はRF信号の増幅・周波数変換のみを行い，gNBは地上局に設置される．このため，地上設備の更新によって世代交代が可能である．一方，Regenerative型では，衛星上にgNB機能（ベースバンド処理，L2/L3制御）が搭載されるため，打ち上げ後にハードウェアを交換することは事実上不可能である．

現在運用中あるいは開発中の多くのNTNプラットフォームは，技術的に成熟した4G LTE（eNB）をベースとしている．地上のネットワークが5G，そして6Gへと進化した際，これら「非地上系ネットワーク上のレガシーRAN」をどのように次世代CNへ収容するかが課題となる．

本ゲートウェイはソフトウェアとして地上側に配置可能であり，衛星やHAPS上のRANを改修することなく，地上のCNのみを5Gへ移行できる可能性を示した．加えて，本ゲートウェイはSCTP接続を終端し，RAN側とCN側のタイムアウト管理を分離する設計となっている．一般的なエンドツーエンドの接続では，長遅延やパケット損失によりSCTPリンク断やNASタイマ満了が発生するケースがあるが，提案ゲートウェイが「プロキシ」として介在しSCTP接続を終端することで，再送制御やタイムアウト値を区間ごとに最適化（例：衛星区間のみタイムアウト値を長く設定するなど）する余地が生まれる．この特性は，静止衛星やHAPSを経由する高遅延なNTN環境において，通信の安定性を向上させるために重要であると考えられる．

NTN地上局が僻地に設置されるケースでは，コア側配置により保守性を確保できる．本研究では高遅延条件（長RTT）を付加した評価や実際のNTN回線（衛星回線等）における実証実験は実施しておらず，これらは今後の課題である．

\section{本章のまとめ}

本章では，提案する「S1-N2変換ゲートウェイ」の機能検証，性能評価，および考察を行った．主要な結果を以下に要約する．

\begin{itemize}
    \item \textbf{アーキテクチャ要件の達成}: R1（透過的介在），R2（プロトコル相互変換），R3（5GC固有機能の整合）の3要件について，PoCによる実証を確認した．ゲートウェイがeNBに対してはMME，5GCに対してはgNBとして振る舞うことで，既存機器に改修を要求しない透過的介在を実現した．
    \item \textbf{設計妥当性の検証}: C-Plane設計において，NAS変換・ID変換・TEIDマッピングが信号レベルで成立することを確認し，「世代共通機能」に基づく設計アプローチの有効性を示した．U-Plane設計において，提案手法とベースライン環境の間に有意な性能差がないことを確認した．
    \item \textbf{疎結合アーキテクチャの意義}: 本アーキテクチャの効果として，(1) RAN/CNのライフサイクル分離（CN更新頻度をRANの制約から切り離せる），(2) プロトコル差分の局所化（検証範囲を境界内に閉じ込められる），(3) 観測点の集約（障害解析やトレースが容易になる）を示した．また，本研究の変換アプローチは短期的には即座に適用可能なゲートウェイとして価値を持ち，長期的にはプロトコル設計への示唆を与えることを論じた．
    \item \textbf{5GC固有機能の活用}: レガシーRANを5GCへ収容することで，スライシング（論理分離），ポリシ制御（QoS高度化），可観測性（クラウドネイティブ運用）などのCN側機能を活用できることを示した．配置戦略としては，コア側配置を基本とし，NTN等の特殊要件がある場合に分散配置を検討する方針が妥当である．
    \item \textbf{応用可能性}: 低コスト検証環境（sXGP＋OSS 5GC），商用マイグレーション（CN先行移行戦略），NTN環境（衛星・HAPS搭載レガシーRANの収容）への適用可能性を示した．特にNTN環境では，ゲートウェイによるSCTP終端が高遅延回線での安定性向上に寄与する可能性がある．
\end{itemize}

次章では，本研究全体の結論として研究課題に対する解決策を整理し，今後の課題と展望を述べる．
